% Generated draft: requires mathematical and source review.
\documentclass[11pt]{article}
\usepackage[margin=25mm]{geometry}
\usepackage{fontspec}
\setmainfont{Latin Modern Roman}
\usepackage{amsmath,amssymb,mathtools,unicode-math}
\setmathfont{Latin Modern Math}
\usepackage{enumitem,needspace,xurl,hyperref}
\setlength{\parindent}{0pt}
\setlength{\parskip}{5pt}
\setlength{\emergencystretch}{3em}
\newcommand{\N}{\mathbb{N}}
\newcommand{\Z}{\mathbb{Z}}
\newcommand{\Q}{\mathbb{Q}}
\newcommand{\R}{\mathbb{R}}
\newcommand{\C}{\mathbb{C}}
\newcommand{\F}{\mathbb{F}}
\newcommand{\A}{\mathbb{A}}
\newcommand{\PP}{\mathbb{P}}
\newcommand{\E}{\mathbb{E}}
\newcommand{\eps}{\varepsilon}
\newcommand{\dd}{\,\mathrm{d}}
\begin{document}
\section*{UnsolvedMath ID: 800934}
\textbf{Source problem code:} catalogue entry
\subsection{Definitions and mathematical statement}
Let \(M\) be a closed connected smooth manifold, where closed means compact and without boundary, and let \(g\) be a Riemannian metric on \(M\). The metric has \emph{no conjugate points} if no nonzero Jacobi field along any geodesic vanishes at two distinct parameter values. Fix a basepoint \(p\in M\) and write \(G=\pi_1(M,p)\).

A finitely generated group \(G\) is \emph{semihyperbolic} if it admits a bounded, equivariant quasigeodesic bicombing. Explicitly, choose a finite symmetric generating set \(S\) and let \(d_S\) be the distance in its Cayley graph. Such a bicombing assigns to each ordered pair \(x,y\in G\) a finite edge path \(\sigma_{x,y}\colon\{0,\ldots,L_{x,y}\}\to G\) from \(x\) to \(y\). Extend this path by \(\sigma_{x,y}(i)=y\) for \(i\ge L_{x,y}\). There must be constants \(\lambda\ge1\) and \(\epsilon,A,B\ge0\), independent of the endpoints, such that
\[
\lambda^{-1}|i-j|-\epsilon
\le d_S\bigl(\sigma_{x,y}(i),\sigma_{x,y}(j)\bigr)
\le \lambda|i-j|+\epsilon
\qquad(0\le i,j\le L_{x,y}),
\]
\[
d_S\bigl(\sigma_{x,y}(i),\sigma_{x',y'}(i)\bigr)
\le A\max\{d_S(x,x'),d_S(y,y')\}+B
\qquad(i\in\mathbb Z_{\ge0}),
\]
and, for every \(h\in G\),
\[
L_{hx,hy}=L_{x,y},\qquad
\sigma_{hx,hy}(i)=h\sigma_{x,y}(i).
\]

\textbf{Question.} For every closed connected Riemannian manifold \((M,g)\) without conjugate points, is \(\pi_1(M,p)\) semihyperbolic?

The Riemannian interpretation and a choice of connected component and basepoint make the fundamental-group question explicit. No dimension restriction appears in the extracted question. The supplied preceding prose mentions an affirmative answer in dimension \(2\), via uniformization, and an unresolved problem even in dimension \(3\); however, its opening reference to an earlier question has no identifiable antecedent in the supplied extract. Those dimensional assertions are therefore preserved as source context without assigning them to this question.
\subsection{Short English statement}
Must every closed Riemannian manifold without conjugate points have a semihyperbolic fundamental group? The registry labels this question open. The extracted dimensional remarks have an unresolved antecedent, so their applicability to this question requires review.
\subsection{Sources}
\begin{itemize}
\item Local UnsolvedMath registry, permanent ID \(800934\), exact source problem code \texttt{catalogue entry}; supplied title: \emph{Open Problems and Questions about Geodesics --- Question 14}. The record attributes the problem to Burns and Matveev and supplies proposed year \(2019\) and status \texttt{open}; these are registry assertions, not independently verified historical or mathematical assessments.
\item Burns and Matveev, \emph{Open Problems and Questions about Geodesics} (source-list year supplied as \(2019\)), Question 14, \url{https://arxiv.org/abs/1308.5417}. The record reports access on 2026-10-16 and extraction from arXiv LaTeX; no independent source visit is claimed here.
\item Source reference labels \texttt{[Cro-Sch1986]}, \texttt{[Lebedeva2002]}, and \texttt{[Iva-Kap2014]}. The supplied text cites these for results comparing properties of fundamental groups of manifolds without conjugate points with those arising in nonpositive curvature. It specifically identifies Section 8 of \texttt{[Iva-Kap2014]} as containing open problems, including the displayed question. Full bibliographic entries and separate links for these references are absent from the supplied record.
\end{itemize}
\end{document}
